\section{XML Data Importing}

\begin{frame}[fragile]{What Is XML Data Importing?}
	\begin{itemize}
		\item XML data files load records when a module is installed or updated.
		\item Perfect for:
		\begin{itemize}
			\item configuration data
			\item sequences
			\item email templates
			\item menus/actions/security
			\item master data that must exist in every database
		\end{itemize}
\begin{block}{Manifest Entry}
\begin{lstlisting}
'data': [
	'data/ir_sequence_data.xml',
	'data/student_data.xml',
],
\end{lstlisting}
\end{block}
	\end{itemize}
\end{frame}

\begin{frame}[fragile]{Basic XML Record Structure}
\begin{lstlisting}[style=xml]
<?xml version="1.0" encoding="utf-8"?>
<odoo>
  <record id="student_demo_ahmed" model="student.student">
	<field name="name">Ahmed Muhumed</field>
	<field name="email">ahmed@example.com</field>
	<field name="age">22</field>
	<field name="gender">male</field>
  </record>
</odoo>
\end{lstlisting}
	\begin{itemize}
		\item \texttt{id}: XML / external ID inside the module.
		\item \texttt{model}: target model technical name.
	\end{itemize}
\end{frame}

\begin{frame}[fragile]{Breaking Down Important XML Features}
\begin{lstlisting}[style=xml]
<!-- Many2one by XML ID -->
<field name="classroom_id" ref="classroom_grade_12"/>

<!-- Boolean / numbers with eval -->
<field name="active" eval="True"/>
<field name="age" eval="18"/>

<!-- Search existing record -->
<field name="country_id" search="[('code', '=', 'SO')]"/>
\end{lstlisting}
	\begin{itemize}
		\item \texttt{ref}: link to another XML ID.
		\item \texttt{eval}: Python expression for non-string values.
		\item \texttt{search}: find a record by domain.
	\end{itemize}
\end{frame}

\begin{frame}[fragile]{noupdate and Forcecreate}
\begin{lstlisting}[style=xml]
<odoo noupdate="1">
  <record id="seq_student_student" model="ir.sequence">
	<field name="name">Student Sequence</field>
	<field name="code">student.student</field>
	<field name="prefix">STU/</field>
	<field name="padding">4</field>
  </record>
</odoo>
\end{lstlisting}
	\begin{itemize}
		\item \texttt{noupdate="1"}: do not overwrite the record on later module updates.
		\item Useful for sequences and data users may customize.
	\end{itemize}
\end{frame}

\begin{frame}[fragile]{Function Calls and Delete in XML}
\begin{lstlisting}[style=xml]
<function model="student.student" name="action_activate_all"/>

<delete model="student.student" search="[('name', '=', 'Temporary')]"/>
\end{lstlisting}
	\begin{itemize}
		\item \texttt{<function>} can call a model method during loading.
		\item \texttt{<delete>} removes matching records.
		\item Use carefully; prefer normal records for most data.
	\end{itemize}
\end{frame}

\begin{frame}[fragile]{XML Data Loading Order}
\begin{lstlisting}
'data': [
	'security/school_security.xml',
	'security/ir.model.access.csv',
	'data/ir_sequence_data.xml',
	'data/classroom_data.xml',
	'data/student_data.xml',
	'views/student_views.xml',
],
\end{lstlisting}
	\begin{itemize}
		\item Load dependencies first (security, parents, referenced records).
		\item A wrong order causes missing XML ID / access errors.
	\end{itemize}
\end{frame}

\begin{frame}{XML Data Best Practices}
	\begin{itemize}
		\item Use stable, descriptive XML IDs.
		\item Keep one concern per file when possible.
		\item Use \texttt{noupdate="1"} for user-customizable config.
		\item Prefer \texttt{ref} over hard-coded database IDs.
		\item Always list files in \texttt{\_\_manifest\_\_.py}.
	\end{itemize}
\end{frame}
